Reinforcement learning (RL) for traffic signal control (TSC) typically relies on simulator-privileged states—such as exact queue lengths, per-vehicle waiting times, and precise network arrivals—that no standard roadside camera can provide, thereby creating a critical deployment gap. To bridge this gap, we constrain both the state representation and the reward function strictly to quantities estimable by a commodity traffic camera utilizing standard bounding-box detection and tracking: a 26-dimensional per-agent observation featuring per-approach aggregation and a signed group-mean phase-imbalance feature, alongside a queue-centric vision-proxy reward that combines a mean-of-squares anti-starvation queue penalty with a signed phase-imbalance pressure term and a local region-of-interest (ROI) exit throughput regularizer. The proposed framework is optimized via parameter-shared Multi-Agent Proximal Policy Optimization (MAPPO) with a centralized per-agent-value-head critic and generalized advantage estimation, explicitly trained within a sublane microsimulation environment under a highly heterogeneous, motorcycle-dominant vehicle mix. Evaluated on a 5-junction corridor and a 16-junction grid under time-varying peak demand, the camera-constrained policy reduces mean travel time by [X]\% compared to fixed-time baselines and [Y]\% versus max-pressure actuation, successfully retaining [Z]\% of the performance achieved by a fully privileged-state policy. Furthermore, under a realistic sensing-noise model with error magnitudes empirically measured from the vision detector stack, the control performance degrades gracefully up to twice the calibrated noise envelope, yielding a rigorously measured information-restriction cost of [W]\% relative to the ideal privileged baseline. Ultimately, quantifying the precise cost of restricting multi-agent control to purely camera-computable signals demonstrates that such deployment is practically viable without simulator-privileged sensing; closed-loop field validation is reported in a companion study.